
\CourseSection{Conclusion}
\label{sec:conclusion}

The ASCC blueprint expands upon the initial proposal by establishing a structured system design, measurable requirements, safety analysis, and detailed development plan. The project aims to integrate golfer following, obstacle detection, emergency stopping, RFID club tracking, and GPS-based course information into a functional autonomous golf cart.

Stakeholder needs have been decomposed into functional, performance, interface, and safety requirements, each linked to planned verification tests. Quantitative targets for following distance, speed, obstacle detection, RFID accuracy, and battery runtime provide measurable acceptance criteria. The hazard analysis further identifies risks involving collisions, uncontrolled movement, electrical faults, and terrain conditions, with corresponding safeguards and residual risks documented.

Development will follow an incremental approach, prioritizing motor control, safety, and tracking before full system integration. Critical functions will be verified through subsystem testing and controlled fault injection before outdoor validation. Technical failures are addressed through defined contingency plans, including alternative tracking methods, RFID configurations, hardware replacements, and adjustments to operating limits. Remaining limitations include sensor blind spots, battery hazards, and restricted terrain suitability.

With an estimated cost of CAD \$480 against the \$600 budget, the project remains financially feasible. The next steps are to finalize quantitative safety thresholds, confirm component compatibility, and complete the initial hardware proofs of concept. Successful completion will be demonstrated through verified subsystem requirements and an integrated prototype capable of supporting a simulated golf hole under controlled conditions.
